% !TEX root = ../main.tex

% 中英标题：\chapter{中文标题}[英文标题]
\chapter{课题的来源及研究的背景和意义}


\section{课题的来源}
蛋白质是生物系统的组成部分，要了解生物系统及其在分子尺度上的行为，就必须了解蛋白质的功能。蛋白质功能预测对于理解细胞内和细胞外生物学过程以及揭示蛋白质多样性与进化之间的关系具有重要意义。此外，蛋白质功能预测在阐明基因功能、理解基因调控网络和指导未表征蛋白质的功能注释方面起着至关重要的作用\cite{eisenberg2000protein}。蛋白质功能的准确预测通过促进靶向特定蛋白质的药物的产生来支持药物发现和开发，从而提高治疗效果。在个性化医疗中，了解蛋白质功能的个体差异对于评估疾病风险和设计个性化治疗策略至关重要。总之，蛋白质功能预测是推动生命科学和医学进步的重要工具，对人类健康和可持续发展具有重要意义。

在生物数据收集技术趋于成熟的今天，人们通过高通量测序和蛋白质组质谱分析技术测定的蛋白质组相关数据呈指数级增长。然而，基于生物实验确定蛋白质功能的方法耗时耗力，并且对设备、资金等方面的要求苛刻，远远不能适应当下日益增长的标注需求。全球最权威的蛋白质数据库 UniProt\cite{uniprot2021uniprot}中只有不到 1$\%$的蛋白质经过实验标注了功能，即便对于研究较广泛的人类蛋白质组已知的2万多个蛋白质而言，也有2/3蛋白质是完全未经实验标注或者功能标签不全的\cite{boutet2016uniprotkb}，亟需进一步标注。在这样的时代背景之下，使用计算手段自动标注蛋白质功能的方法应运而生，其相比传统的实验标注方法更加方便高效、预测成本低、扩展性强且应用范围广，因此逐渐成为蛋白质功能标注领域的研究热点。

\section{课题研究的背景和意义}
蛋白质序列中的氨基酸残基通过一系列复杂的空间转化过程形成独特的三维结构，这种结构在生物学功能中起着至关重要的作用。蛋白质的三维结构不仅赋予了其独特的物理化学性质，还通过非共价相互作用形成蛋白质-蛋白质相互作用（Protein-Protein Interaction，PPI）网络。这些网络是细胞中生命活动的核心调控机制，参与了几乎所有的生物过程。例如PPI网络可以通过组装成蛋白质复合物，形成特定的分子机器，来执行复杂的生物学功能。然而，值得注意的是并非所有蛋白质都需要通过形成复合物来实现其功能，有些蛋白质能够独立地执行其功能，这使得蛋白质功能的多样性和复杂性更加显著。

在蛋白质研究领域，遵循“序列-结构-功能”的范式，即通过蛋白质的氨基酸序列推测其三维结构，再通过结构推测其生物功能。目前，蛋白质功能的注释和分类主要依赖于基因本体论（Gene Ontology，GO）系统。GO系统是一个标准化的框架，涵盖了三大本体论：生物过程本体论（Biological Process Ontology，BPO）、分子功能本体论（Molecular Function Ontology，MFO）和细胞组分本体论（Cellular Component Ontology，CCO）。GO系统目前包含超过50,000个类目，这些类目之间并非独立存在，而是通过正式定义的关系相互连接和影响。因此，在预测蛋白质功能时，必须考虑这些类目之间的复杂关系。尽管GO系统为蛋白质功能的注释提供了全面的框架，仍有许多GO类目由于缺乏足够的实验注释数据而难以准确预测。事实上，在GO系统的20,000多个类目中，有超过一半的类目仅有极少数的蛋白质被实验注释为具有该功能，甚至一些类目完全没有注释数据。这些稀缺注释类别通常是高度特异的，具有较高的信息量，对生物学研究和应用具有重要意义。由于数据稀缺，传统的机器学习模型在这些类别上的预测表现往往较差，随着越来越多的新蛋白质被发现，如何准确预测这些新蛋白质的功能成为了蛋白质研究领域的一个重大挑战。

在这种背景下，基于多模态预训练的方法为解决这一问题提供了一种创新的解决方案\cite{wang2023large}。多模态预训练方法旨在通过整合来自蛋白质的不同数据类型，包括序列、结构、以及可能的功能注释等来训练预测模型，能够更全面地捕捉蛋白质的复杂性。相较于传统的依赖大量标注数据的方法，多模态预训练不仅可以处理标注数据稀缺的问题，还能提高模型对新蛋白质功能的预测精度。本研究的目的在于开发和验证一种基于多模态预训练的蛋白质功能预测模型，通过这一模型，不仅能够有效地预测那些未标注蛋白质的功能，还能够扩展现有的蛋白质功能预测方法的应用范围，提升其准确性和广泛性。这将有助于推动生命科学的进步，为疾病治疗和药物开发等应用领域提供强有力的科学基础。